% ============================================================
% Autor: Eloi Egea Rada
% Tema: Existing solutions, platforms, and gap analysis
% ============================================================

\chapter{State of the Art}
\label{ch:state-of-art}

Cybersecurity education has experienced exponential growth over the past decade, driven by
increasing industry demand for skilled professionals and the accessibility of online learning
platforms. While students have unprecedented access to training resources, the fragmentation
of educational approaches has created gaps between abstract theoretical knowledge and practical,
integrated system security. This chapter examines existing platforms and methodologies, identifies
gaps in current offerings, and positions this bachelor's thesis within the broader context of
cybersecurity pedagogical innovation.


\section{Existing Platforms Overview}

The major platforms used for cybersecurity training share a common limitation: they are designed
for independent learners or certification preparation rather than curriculum-aligned institutional
deployment. Table~\ref{tab:platform_comparison} summarises their key characteristics.
Platforms such as HackTheBox~\cite{hackthebox} and TryHackMe~\cite{tryhackme} are externally hosted
and cannot be deployed on institutional infrastructure. SEED Labs~\cite{seedlabs} and
GOAD~\cite{goad} support local deployment but cover only isolated concepts or narrow domains.
VulnHub~\cite{vulnhub} and Hack4u~\cite{hack4u} lack automation and assessment integration.
Cisco Networking Academy~\cite{cisco_netacad} provides formal structure but is tightly coupled
to vendor technology and limited to network simulation.

\begin{table}[H]
\centering
\scriptsize
\setlength{\tabcolsep}{4pt}
\begin{tabularx}{\textwidth}{lXXXXX}
\toprule
\textbf{Platform} & \textbf{Approach} & \textbf{Deploy} & \textbf{Curriculum} & \textbf{Auto} & \textbf{Network} \\
\midrule
HackTheBox & CTF & No & Minimal & No & Isolated \\
TryHackMe & Gamified & No & Structured & No & Limited \\
Cisco Academy & Academic & Yes (PT) & Strong & Limited & Net-focused \\
SEED Labs & Conceptual & Yes (Docker) & Concept-based & Minimal & Isolated \\
GOAD & AD-focused & Yes (\mbox{Terraform}) & None & Moderate & \mbox{AD Domain} \\
VulnHub & Community & Yes & None & No & Isolated \\
Hack4u Academy & Methodology & External & Strong & Minimal & Limited \\
\bottomrule
\end{tabularx}
\caption{Comparative Analysis: Existing Cybersecurity Training Platforms}
\label{tab:platform_comparison}
\end{table}

Full individual platform descriptions, including detailed discussion of each platform's
pedagogical approach, deployment model, and specific limitations, are provided in
Appendix~\ref{app:platforms}.

\section{Security Frameworks and Methodologies}

Several authoritative frameworks inform cybersecurity education and guide both attack
methodologies and defensive best practices.

\subsection{OWASP Top 10}

The OWASP Top 10 is a regularly updated list of the most critical web application security
vulnerabilities \cite{owasp2021}. For educational purposes it provides a focused, industry-recognised
curriculum that prioritises the attack vectors students are most likely to encounter in practice.
Lab 2 of this project is directly aligned with OWASP Top 10 categories.

\subsection{CIS Critical Controls}

The CIS Critical Controls (v8.0) represent a prioritised set of defensive techniques spanning
network defence, data protection, and incident response \cite{cis_controls}. Labs 1 and 3
incorporate defensive controls so that students understand not only how to attack but why
specific security measures are implemented.

\subsection{NIST Cybersecurity Framework}

The NIST Cybersecurity Framework organises security activities into five functions: Identify,
Protect, Detect, Respond, and Recover \cite{nist_framework}. The lab design follows these
principles by enabling students to identify vulnerabilities through hands-on exercises and
understand corresponding protective measures.




\section{Prior Academic Research}

Academic research consistently validates the effectiveness of hands-on, laboratory-based
approaches in cybersecurity education. Conti and Di Pietro \cite{conti2019} demonstrate that
challenge-based environments significantly improve skill retention and practical competency
compared to purely theoretical instruction. Murphy et al.\ \cite{murphy2018} provide a
structured methodology for designing laboratory exercises that are explicitly aligned with
curriculum outcomes, validating the need for scenario-specific tasks and integrated
assessment. Sponj and Heri\v{c}ko \cite{sponj2021} conducted a systematic review of
cybersecurity education research, identifying hands-on practice and authentic task design as
the two most influential factors in learner engagement and competency development.
Al~Mohanadi and Furnell \cite{almohannadi2016} further establish a requirements-driven
approach to information security assessment laboratories, emphasising that lab tasks must
map directly to defined learning outcomes to achieve meaningful pedagogical impact.

These findings directly inform the design choices in this project: each lab is built around
scenario-based flags mapped to specific learning objectives, includes an instructor reference
sheet, and culminates in an automated assessment mechanism, addressing the gap between
platform availability and curriculum integration identified in the prior literature.

\section{Gap Analysis: What Existing Solutions Miss}

Despite the proliferation of cybersecurity training platforms, significant gaps remain:

\textbf{Gap 1: Curriculum Integration.}
Most platforms target independent learners or certification preparation rather than integration
into specific degree programmes. No major platform explicitly maps pentesting scenarios to the
learning outcomes of a particular computer science curriculum.

\textbf{Gap 2: Institution Control.}
Effective educational deployment requires institutional control over content, assessment, and
infrastructure. External platforms introduce dependencies and limit customisation, while
fully self-hosted solutions often impose high maintenance burdens.

\textbf{Gap 3: Integrated Network Penetration Testing.}
Most platforms focus on isolated services or individual techniques. Few provide realistic,
multi-system network topologies where reconnaissance informs a complete attack chain and
lateral movement is required.

\textbf{Gap 4: Automation and Scalability.}
Even classroom-oriented platforms often lack automated deployment, environment reset, and
assessment mechanisms, increasing instructor workload and limiting large-scale institutional
adoption.

\textbf{Gap 5: Balance of Guidance and Autonomy.}
Effective labs must balance sufficient structure to keep students on target with sufficient
openness to encourage problem-solving and methodology development.


\section{The Need for EVE-NG Based Labs}

This analysis reveals that while excellent individual components exist---gamified learning
(TryHackMe), challenging scenarios (HackTheBox), research-backed content (SEED Labs),
specialised domains (GOAD)---no existing solution comprehensively addresses the needs of a
university cybersecurity education programme:

\begin{itemize}
    \item \textbf{Curriculum-aligned labs}: Mapped to specific degree learning outcomes
    \item \textbf{Integrated network scenarios}: Multi-system penetration testing with realistic complexity
    \item \textbf{Institutional control}: Locally deployable, customisable, and maintainable
    \item \textbf{Complete automation}: One-command deployment and reset with built-in assessment
    \item \textbf{System-level realism}: Full virtual machine workflows (real OS images, ISO-based
          installation, native kernel behaviour) that reproduce the technical conditions of
          professional pentesting practice
\end{itemize}

The bachelor's thesis project bridges this gap by leveraging EVE-NG \cite{eve_ng} to create a
comprehensive, reusable, and automatable teaching package for cybersecurity practical labs,
synthesising the pedagogical strengths of existing platforms while addressing their limitations
through institutional control, curriculum alignment, and automation.


% END OF CHAPTER 2: STATE OF THE ART
